\chapter{Conclusão}
\label{chap:conclusao}

Este trabalho investigou como a eficácia de estratégias de retreinamento
varia em função do tipo e da severidade do drift e do classificador
monitorado, por meio de um experimento fatorial controlado com 180 condições
e 1.800 execuções em ambiente de stream. A severidade foi operacionalizada
pela distância de Hellinger, cuja forma fechada entre distribuições normais
permitiu injetar cada nível com valor exato e compará-lo com a estimativa do
detector.

Os resultados mostram que tipo, severidade e classificador interagem
significativamente com a estratégia, mas que, nas condições estudadas, o que
mais diferenciou as estratégias foi a forma de acionamento. A janela
deslizante, retreinada em cadência fixa, foi a mais eficaz em 29 das 36
combinações de tipo, severidade e classificador, com vantagem pequena sob
drift abrupto e recorrente e grande sob drift gradual e incremental, nos
quais o HDDDM detecta a mudança no meio da transição e as estratégias
acionadas por ele retreinam cedo demais. Essa vantagem custou cerca de quatro
vezes mais retreinamentos e desapareceu quando a cadência foi reduzida. O
reset completo foi a melhor estratégia sob drift abrupto de severidade alta;
a reativação de modelos históricos identificou corretamente as recorrências,
mas igualou o reset em acurácia, com menor custo; e o retreinamento
incremental dependeu fortemente do classificador, sendo equivalente ao
reset no XGBoost e o pior entre as estratégias no Random Forest.

Quanto à questão derivada, os sinais observáveis do detector mostraram-se
informativos: a severidade observada acompanhou a severidade real, a
velocidade da divergência distinguiu mudanças abruptas de lentas e a
distância ao repositório identificou as recorrências. A orientação construída
a partir deles superou a prática de sempre aplicar o reset completo, mas não
a política fixa de sempre usar a janela deslizante, já que esta foi quase
sempre a estratégia mais eficaz.

As principais contribuições do trabalho são: (i) um protocolo experimental
reprodutível, disponibilizado em um notebook executável, que controla tipo e
severidade do drift por construção e permite comparar a severidade real e a
observada; (ii) evidência de que, para drift lento, o momento do acionamento
importa mais que o procedimento de retreinamento, de modo que estratégias
acionadas por detecção podem perder para uma janela agendada mesmo usando o
mesmo procedimento de ajuste; (iii) evidência de que a eficácia da
atualização incremental depende do mecanismo do classificador; e (iv) a
indicação de que a estatística de um detector baseado em distribuição, sem
uso de rótulos, carrega informação suficiente para estimar severidade,
velocidade e recorrência do drift.

\section{Trabalhos Futuros}

Os limites discutidos na Seção~\ref{sec:ameacas} apontam os desdobramentos
mais diretos. O primeiro é replicar o experimento com dados reais e com os
conjuntos de referência usados pelo estudo de base \cite{rocha2026}, com
atributos correlacionados e transições de duração variável. O segundo é
incluir os demais detectores baseados em distribuição (JSDDM e KSDDM) e
estratégias híbridas, que combinem acionamento por detecção com
retreinamento agendado ou que continuem retreinando após uma detecção até a
estabilização da divergência. O terceiro é construir uma orientação sensível
ao custo, que pondere acurácia e número de retreinamentos e use os sinais
observáveis para decidir quando uma estratégia mais barata é suficiente,
avaliada como política online, que troque de estratégia a cada episódio. Por
fim, cenários com janelas de treino maiores e rótulos atrasados permitiriam
verificar se a reativação de modelos históricos, que neste trabalho igualou
o reset, passa a superá-lo quando o modelo reativado foi treinado com mais
dados do que os disponíveis após a mudança.
